\documentclass[12pt,a4paper]{article}

\usepackage[utf8]{inputenc}
\usepackage[T1]{fontenc}
\usepackage[spanish]{babel}
\usepackage{geometry}
\usepackage{hyperref}
\usepackage{listings}
\usepackage{xcolor}
\usepackage{enumitem}
\usepackage{parskip}
\usepackage{titlesec}
\usepackage{booktabs}
\usepackage{array}

\geometry{
  top=2.5cm,
  bottom=2.5cm,
  left=3cm,
  right=3cm
}

% Estilo para bloques de código/comandos
\definecolor{codegray}{rgb}{0.95,0.95,0.95}
\definecolor{codegreen}{rgb}{0.0,0.45,0.0}
\definecolor{codepurple}{rgb}{0.45,0.0,0.55}
\definecolor{codeblue}{rgb}{0.0,0.0,0.6}

\lstdefinestyle{bash}{
  backgroundcolor=\color{codegray},
  basicstyle=\ttfamily\small,
  keywordstyle=\color{codeblue},
  commentstyle=\color{codegreen},
  stringstyle=\color{codepurple},
  breakatwhitespace=false,
  breaklines=true,
  captionpos=b,
  keepspaces=true,
  showspaces=false,
  showstringspaces=false,
  frame=single,
  rulecolor=\color{gray!40},
  xleftmargin=0.5em,
  xrightmargin=0.5em,
  aboveskip=0.8em,
  belowskip=0.8em
}

\lstdefinestyle{config}{
  backgroundcolor=\color{codegray},
  basicstyle=\ttfamily\small,
  breaklines=true,
  frame=single,
  rulecolor=\color{gray!40},
  xleftmargin=0.5em,
  xrightmargin=0.5em,
  aboveskip=0.8em,
  belowskip=0.8em
}

\hypersetup{
  colorlinks=true,
  linkcolor=codeblue,
  urlcolor=codeblue
}

\setlength{\parskip}{0.6em}

\title{\textbf{Configuración de MQTT con TLS en Raspberry Pi 4}\\
\large Habilitación exclusiva del puerto 8883}
\author{}
\date{}

\begin{document}

\maketitle
\thispagestyle{empty}

% ============================================================
\section{¿Qué es MQTT?}
% ============================================================

MQTT (del inglés \textit{Message Queuing Telemetry Transport}) es un protocolo de
comunicación diseñado para transmitir mensajes de forma ligera y eficiente entre
dispositivos. Fue creado originalmente para situaciones donde la conexión a
internet es lenta, inestable o donde los dispositivos tienen recursos muy
limitados, como microcontroladores con poca memoria.

Hoy en día, MQTT se utiliza ampliamente en el \textbf{Internet de las Cosas}
(\textit{IoT}, por sus siglas en inglés), que es el conjunto de dispositivos
físicos conectados a internet para enviar y recibir información: sensores de
temperatura, cámaras, actuadores, medidores de energía, etc.

\subsection{Modelo de publicación y suscripción}

MQTT funciona con un modelo llamado \textit{publicación/suscripción}
(``pub/sub''). A diferencia de una llamada telefónica directa entre dos
personas, este modelo introduce un intermediario llamado \textbf{bróker}. El
funcionamiento general es el siguiente:

\begin{itemize}
  \item Un dispositivo que tiene datos (por ejemplo, un sensor de temperatura)
    los \textbf{publica} en el bróker bajo un nombre llamado \textbf{tópico}
    (en inglés, \textit{topic}). Por ejemplo: \texttt{casa/sala/temperatura}.
  \item Otro dispositivo que quiere recibir esos datos se \textbf{suscribe} al
    mismo tópico en el bróker.
  \item El bróker actúa como punto central: recibe los mensajes y los reenvía a
    todos los suscriptores del tópico correspondiente.
\end{itemize}

Con este esquema, el dispositivo que publica no necesita saber quién va a
recibir el mensaje, y el dispositivo que recibe no necesita saber de dónde
proviene. Todo pasa a través del bróker.

\subsection{Ventajas de MQTT en IoT}

\begin{itemize}
  \item \textbf{Bajo consumo de ancho de banda:} los mensajes son muy pequeños,
    ideales para dispositivos con conexión lenta o de bajo costo.
  \item \textbf{Bajo consumo energético:} pensado para dispositivos con batería.
  \item \textbf{Simplicidad:} fácil de implementar en hardware con recursos
    limitados.
  \item \textbf{Flexibilidad:} un bróker puede manejar miles de clientes al
    mismo tiempo.
\end{itemize}

% ============================================================
\section{El bróker Mosquitto en Raspberry Pi 4}
% ============================================================

El \textbf{bróker} es el servidor central de MQTT: todos los dispositivos se
conectan a él para publicar o suscribirse. En este proyecto se utiliza
\textbf{Eclipse Mosquitto}, una implementación de código abierto, muy liviana y
ampliamente utilizada en proyectos IoT.

La \textbf{Raspberry Pi 4} funciona como el servidor (bróker) donde se instala y
ejecuta Mosquitto. Es un mini-ordenador de bajo costo y consumo que corre el
sistema operativo Raspberry Pi OS (basado en Linux/Debian), lo que lo hace ideal
para este tipo de aplicaciones.

\subsection{Puertos de comunicación}

MQTT utiliza puertos TCP para las conexiones:

\begin{itemize}
  \item \textbf{Puerto 1883:} conexión sin cifrado. Cualquier persona que pueda
    interceptar el tráfico de red podría leer los mensajes.
  \item \textbf{Puerto 8883:} conexión cifrada mediante \textbf{TLS/SSL}
    (el mismo mecanismo que usa HTTPS en los navegadores web). Garantiza que
    los mensajes viajan de forma segura y que el bróker es quien dice ser.
\end{itemize}

En este documento se configurará el bróker para que acepte \textbf{únicamente}
conexiones en el puerto 8883, es decir, solo conexiones seguras y cifradas.

\subsection{¿Qué es TLS y por qué es importante?}

TLS (\textit{Transport Layer Security}) es un protocolo de seguridad que cifra
la comunicación entre dos partes. Para que TLS funcione se necesitan
\textbf{certificados digitales}, que son archivos que identifican a un servidor
de forma confiable, similar a un documento de identidad digital.

Para generar estos certificados se necesita una \textbf{Autoridad Certificadora}
(CA, del inglés \textit{Certificate Authority}), que es quien ``firma'' y avala
los certificados. En este caso se creará una CA propia (autofirmada), válida
para uso interno o en redes locales.

% ============================================================
\section{Configuración del bróker Mosquitto con TLS\\(puerto 8883 exclusivo)}
% ============================================================

A continuación se describe el procedimiento completo para instalar Mosquitto en
una Raspberry Pi 4 y configurarlo para operar \textbf{exclusivamente} en el
puerto 8883 con TLS.

Todos los comandos se ejecutan en la terminal de la Raspberry Pi. Los que
comienzan con \texttt{sudo} requieren permisos de administrador.

% ----------------------------------------------------------
\subsection{Paso 1: Limpieza e instalación de Mosquitto}
% ----------------------------------------------------------

Si el sistema ya tenía una instalación previa de Mosquitto (por ejemplo,
configurada solo para el puerto 1883), es recomendable eliminarla
completamente para evitar conflictos con archivos de configuración anteriores.

\begin{lstlisting}[style=bash]
sudo systemctl stop mosquitto 2>/dev/null || true
sudo apt purge -y mosquitto mosquitto-clients
sudo rm -rf /etc/mosquitto /var/lib/mosquitto /var/log/mosquitto
\end{lstlisting}

Tras limpiar el sistema, se actualiza la lista de paquetes disponibles y se
instala Mosquitto junto con sus herramientas de cliente:

\begin{lstlisting}[style=bash]
sudo apt update
sudo apt install -y mosquitto mosquitto-clients
\end{lstlisting}

Para verificar que la instalación fue exitosa:

\begin{lstlisting}[style=bash]
mosquitto -h | head -1
systemctl status mosquitto --no-pager -l
\end{lstlisting}

% ----------------------------------------------------------
\subsection{Paso 2: Crear la estructura de directorios y usuarios}
% ----------------------------------------------------------

Mosquitto organiza su configuración en carpetas específicas. Se crean los
directorios necesarios con el siguiente comando:

\begin{lstlisting}[style=bash]
sudo mkdir -p /etc/mosquitto/{conf.d,certs,ca_certificates,passwords}
\end{lstlisting}

Cada directorio cumple una función concreta:

\begin{itemize}
  \item \texttt{conf.d/}: archivos de configuración adicionales que Mosquitto
    carga automáticamente al iniciarse.
  \item \texttt{certs/}: almacena la clave privada y el certificado del bróker.
  \item \texttt{ca\_certificates/}: almacena el certificado de la CA
    que validará las conexiones de los clientes.
  \item \texttt{passwords/}: almacena el archivo de usuarios y contraseñas.
\end{itemize}

A continuación se crea el usuario \texttt{usuario1} con contraseña
\texttt{qwerty123}. El comando \texttt{mosquitto\_passwd} guarda la contraseña
de forma \textbf{cifrada} usando un algoritmo de hash seguro (PBKDF2 con
SHA-512 en versiones recientes de Mosquitto), lo que la hace resistente a
ataques de diccionario:

\begin{lstlisting}[style=bash]
sudo mosquitto_passwd -c /etc/mosquitto/passwords/users.db usuario1
\end{lstlisting}

El sistema pedirá escribir la contraseña dos veces para confirmarla.

Luego se asignan los permisos correctos al archivo de contraseñas:

\begin{lstlisting}[style=bash]
sudo chown root:mosquitto /etc/mosquitto/passwords/users.db
sudo chmod 640 /etc/mosquitto/passwords/users.db
\end{lstlisting}

\textbf{¿Por qué estos permisos?}

\begin{itemize}
  \item \texttt{chown root:mosquitto} asigna como dueño a \texttt{root} (el
    administrador del sistema) y como grupo a \texttt{mosquitto} (el usuario
    con el que corre el servicio).
  \item \texttt{chmod 640} establece los permisos en octal:
    \begin{itemize}
      \item \textbf{6 (rw-)} para el dueño: \texttt{root} puede leer y escribir.
      \item \textbf{4 (r--)} para el grupo: \texttt{mosquitto} solo puede leer
        (necesario para autenticar a los clientes).
      \item \textbf{0 (---)} para todos los demás: nadie más puede acceder al
        archivo.
    \end{itemize}
\end{itemize}

En resumen: \texttt{root} controla el archivo, el servicio Mosquitto puede
leerlo para autenticar conexiones, y ningún otro usuario del sistema tiene
acceso.

% ----------------------------------------------------------
\subsection{Paso 3: Generar los certificados TLS}
% ----------------------------------------------------------

Para cifrar las comunicaciones en el puerto 8883 se necesitan tres elementos:

\begin{enumerate}
  \item \textbf{Certificado de la CA (ca.crt):} es el certificado raíz que
    firma y valida los demás certificados.
  \item \textbf{Clave privada del bróker (broker.key):} archivo secreto que
    solo debe conocer el bróker.
  \item \textbf{Certificado del bróker (broker.crt):} identifica al bróker y
    está firmado por la CA.
\end{enumerate}

\subsubsection{3.1 Crear la Autoridad Certificadora (CA)}

\begin{lstlisting}[style=bash]
# Generar la clave privada de la CA (2048 bits)
openssl genrsa -out /etc/mosquitto/ca_certificates/ca.key 2048

# Generar el certificado autofirmado de la CA (válido por 3650 días ~10 años)
openssl req -new -x509 -days 3650 \
  -key /etc/mosquitto/ca_certificates/ca.key \
  -out /etc/mosquitto/ca_certificates/ca.crt \
  -subj "/CN=MiCA-MQTT"
\end{lstlisting}

\subsubsection{3.2 Crear la clave y el certificado del bróker}

\begin{lstlisting}[style=bash]
# Generar la clave privada del broker
openssl genrsa -out /etc/mosquitto/certs/broker.key 2048

# Crear la solicitud de certificado (CSR)
openssl req -new \
  -key /etc/mosquitto/certs/broker.key \
  -out /etc/mosquitto/certs/broker.csr \
  -subj "/CN=broker-mqtt"
\end{lstlisting}

\subsubsection{3.3 Agregar la dirección IP como SAN (\textit{Subject Alternative Name})}

El campo SAN (\textit{Subject Alternative Name}) indica a qué direcciones IP o
nombres de dominio es válido el certificado. Esto es \textbf{obligatorio} en
las versiones modernas de TLS: sin SAN, los clientes rechazarán el certificado
con un error de verificación.

Se crea un archivo temporal con la extensión de SAN, indicando la IP de la
Raspberry Pi (en este ejemplo, \texttt{192.168.1.100}; se debe reemplazar con
la IP real del dispositivo):

\begin{lstlisting}[style=bash]
cat > /tmp/broker-ext.cnf << EOF
[v3_req]
subjectAltName = IP:192.168.1.100
EOF
\end{lstlisting}

\subsubsection{3.4 Firmar el certificado del bróker con la CA}

\begin{lstlisting}[style=bash]
openssl x509 -req -days 3650 \
  -in /etc/mosquitto/certs/broker.csr \
  -CA /etc/mosquitto/ca_certificates/ca.crt \
  -CAkey /etc/mosquitto/ca_certificates/ca.key \
  -CAcreateserial \
  -out /etc/mosquitto/certs/broker.crt \
  -extensions v3_req \
  -extfile /tmp/broker-ext.cnf
\end{lstlisting}

\textbf{Nota sobre la validez del certificado:} el valor \texttt{-days 3650}
establece una vigencia de 10 años, suficiente para entornos de desarrollo o
redes internas. En un entorno de producción se recomienda usar períodos más
cortos (por ejemplo, 365 días) y renovar los certificados periódicamente.

\subsubsection{3.5 Ajustar los permisos de los certificados}

\begin{lstlisting}[style=bash]
sudo chown root:mosquitto /etc/mosquitto/certs/broker.key
sudo chmod 640 /etc/mosquitto/certs/broker.key
sudo chmod 644 /etc/mosquitto/certs/broker.crt
sudo chmod 644 /etc/mosquitto/ca_certificates/ca.crt
\end{lstlisting}

La clave privada (\texttt{broker.key}) es el archivo más sensible: solo
\texttt{root} puede modificarla y solo el grupo \texttt{mosquitto} puede
leerla. Los certificados públicos (\texttt{.crt}) pueden ser leídos por
cualquier usuario del sistema, ya que no contienen información secreta.

\subsubsection{3.6 Verificar los certificados}

Para confirmar que el certificado del bróker fue firmado correctamente por la
CA y que el SAN está incluido:

\begin{lstlisting}[style=bash]
openssl verify -CAfile /etc/mosquitto/ca_certificates/ca.crt \
  /etc/mosquitto/certs/broker.crt

openssl x509 -in /etc/mosquitto/certs/broker.crt -noout -text \
  | grep -A2 "Subject Alternative Name"
\end{lstlisting}

El primer comando debe responder \texttt{OK}. El segundo debe mostrar la IP
configurada en el paso anterior.

% ----------------------------------------------------------
\subsection{Paso 4: Configurar Mosquitto para usar solo el puerto 8883}
% ----------------------------------------------------------

Se crea el archivo de configuración principal en el directorio
\texttt{conf.d/}. Mosquitto carga automáticamente todos los archivos
\texttt{.conf} que encuentre en esa carpeta al iniciar.

\begin{lstlisting}[style=bash]
sudo nano /etc/mosquitto/conf.d/mqtt_8883.conf
\end{lstlisting}

El contenido del archivo debe ser el siguiente:

\begin{lstlisting}[style=config]
# Puerto seguro con TLS/SSL
listener 8883

# Certificados TLS
cafile   /etc/mosquitto/ca_certificates/ca.crt
certfile /etc/mosquitto/certs/broker.crt
keyfile  /etc/mosquitto/certs/broker.key

# Versión mínima de TLS (recomendado TLS 1.2 o superior)
tls_version tlsv1.2

# Requerir usuario y contraseña para conectarse
allow_anonymous false
password_file /etc/mosquitto/passwords/users.db
\end{lstlisting}

\textbf{Descripción de cada directiva:}

\begin{itemize}
  \item \texttt{listener 8883}: indica a Mosquitto que escuche conexiones
    entrantes únicamente en el puerto 8883. Al no agregar un bloque
    \texttt{listener 1883}, ese puerto queda \textbf{deshabilitado}.
  \item \texttt{cafile}: ruta al certificado de la CA. Los clientes deben
    presentar un certificado firmado por esta CA (o simplemente tener la CA en
    su almacén de confianza).
  \item \texttt{certfile} y \texttt{keyfile}: certificado público y clave
    privada del bróker, usados para establecer el cifrado TLS.
  \item \texttt{tls\_version tlsv1.2}: exige al menos TLS 1.2, descartando
    versiones antiguas e inseguras (SSLv3, TLS 1.0, TLS 1.1).
  \item \texttt{allow\_anonymous false}: prohíbe conexiones sin autenticación.
    Todo cliente debe identificarse con usuario y contraseña.
  \item \texttt{password\_file}: apunta al archivo de usuarios y contraseñas
    creado en el Paso 2.
\end{itemize}

% ----------------------------------------------------------
\subsection{Paso 5: Reiniciar y verificar el servicio}
% ----------------------------------------------------------

Se reinicia Mosquitto para que cargue la nueva configuración:

\begin{lstlisting}[style=bash]
sudo systemctl restart mosquitto
sudo systemctl status mosquitto --no-pager -l
\end{lstlisting}

El estado debe aparecer como \texttt{active (running)}. Para confirmar que
Mosquitto está escuchando \textbf{solo} en el puerto 8883 y no en el 1883:

\begin{lstlisting}[style=bash]
sudo ss -tlnp | grep mosquitto
\end{lstlisting}

La salida debe mostrar únicamente el puerto \textbf{8883} asociado al proceso
\texttt{mosquitto}.

% ----------------------------------------------------------
\subsection{Paso 6: Prueba de conectividad}
% ----------------------------------------------------------

Para verificar que el bróker funciona correctamente se utilizan las
herramientas de cliente que incluye el paquete \texttt{mosquitto-clients}.

Abrir dos terminales en la Raspberry Pi (o en cualquier equipo de la red con
acceso al bróker):

\textbf{Terminal A — Suscriptor} (espera mensajes en el tópico
\texttt{prueba/sensor}):

\begin{lstlisting}[style=bash]
mosquitto_sub -h 192.168.1.100 -p 8883 \
  --cafile /etc/mosquitto/ca_certificates/ca.crt \
  -u usuario1 -P qwerty123 \
  -t prueba/sensor -v
\end{lstlisting}

\textbf{Terminal B — Publicador} (envía un mensaje al mismo tópico):

\begin{lstlisting}[style=bash]
mosquitto_pub -h 192.168.1.100 -p 8883 \
  --cafile /etc/mosquitto/ca_certificates/ca.crt \
  -u usuario1 -P qwerty123 \
  -t prueba/sensor -m "Hola desde MQTT con TLS"
\end{lstlisting}

Si la configuración es correcta, la Terminal A mostrará el mensaje:

\begin{lstlisting}[style=config]
prueba/sensor Hola desde MQTT con TLS
\end{lstlisting}

\textbf{Explicación de los parámetros usados:}

\begin{itemize}
  \item \texttt{-h}: dirección IP del bróker.
  \item \texttt{-p 8883}: puerto de conexión.
  \item \texttt{--cafile}: certificado de la CA para verificar la identidad del
    bróker y establecer el cifrado TLS.
  \item \texttt{-u} y \texttt{-P}: usuario y contraseña de autenticación.
  \item \texttt{-t}: nombre del tópico.
  \item \texttt{-m}: mensaje a publicar (solo en el publicador).
  \item \texttt{-v}: modo detallado, muestra el tópico junto al mensaje (solo
    en el suscriptor).
\end{itemize}

% ============================================================
\section{Resumen}
% ============================================================

En este documento se describe cómo funciona MQTT en el contexto del IoT y cómo
configurar un bróker Mosquitto en una Raspberry Pi 4 para que opere
\textbf{exclusivamente} con conexiones seguras en el puerto 8883.

Los puntos clave del proceso son:

\begin{enumerate}
  \item Instalar Mosquitto desde los repositorios oficiales de Raspberry Pi OS.
  \item Crear la estructura de directorios y un archivo de usuarios con
    contraseña cifrada.
  \item Generar una Autoridad Certificadora propia y firmar un certificado para
    el bróker con el campo SAN apuntando a la IP del dispositivo.
  \item Configurar Mosquitto para escuchar \textbf{solo} en el puerto 8883 con
    TLS 1.2 y autenticación obligatoria.
  \item Verificar que el servicio esté activo y comprobar la comunicación con
    las herramientas \texttt{mosquitto\_pub} y \texttt{mosquitto\_sub}.
\end{enumerate}

Con esta configuración, el bróker MQTT queda protegido: todas las
comunicaciones viajan cifradas (no pueden ser leídas por terceros en la red) y
solo los usuarios autorizados pueden conectarse al sistema.

\end{document}
